\subsection{\texorpdfstring{Properties of homomorphisms \(\pi_{1}(f,a)\)}{Properties of homomorphisms pi\_1(f,a)}}

Let A and B be topological spaces, let \((\mathrm{E},p)\) be a covering of A, let \((\mathrm{E}^{\prime},p^{\prime})\) be a covering of B, let \(f\colon A\to B\) and \(g\colon E\to E^{\prime}\) be continuous maps such that \(p^{\prime}\circ g=f\circ p\) . Let a be a point of A and let \(b=f(a)\) . For all \(\gamma\in\pi_{1}(\mathrm{A},a)\) and every \(x\in E_{a}\) , we have \(g(x\cdot\gamma)=g(x)\cdot f_{*}(\gamma)\) , according to relation (3), III, p. 304. If the diagram

\[
\begin{array}{c} \mathrm{E} \xrightarrow {g} \mathrm{E} ^ {\prime} \\ \Big \downarrow p \\ \mathrm{A} \xrightarrow {f} \mathrm{B} \end{array}
\]

is a Cartesian square, the map g induces a bijection from \(E_{a}\) onto \(E'_{b}\) . The action of \(\pi_{1}(A,a)\) on the fiber \(E_{a}\) is then the composite of the action of the group \(\pi_{1}(B,b)\) on \(E'_{b}\) and the homomorphism \(\pi_{1}(f,a)\) from \(\pi_{1}(A,a)\) into \(\pi_{1}(B,b)\) .

PROPOSITION 1. --- Let A be a loop-disentanglable topological space, B a locally path-connected topological space, and let \(f\colon\mathrm{A}\to\mathrm{B}\) be a continuous map. Let a be a point of A and \(b=f(a)\) . Suppose that every covering of A is isomorphic to the inverse image under \(f\) of a covering of B. Then the homomorphism \(\pi_{1}(f,a)\colon\pi_{1}(\mathrm{A},a)\to\pi_{1}(\mathrm{B},b)\) is injective.

Since the space A is loop-disentanglable, there exists a covering E of A whose fiber \(E_{a}\) is a principal homogeneous \(\pi_{1}(A,a)\) -set (IV, p. 342, theorem 1). By hypothesis, this action is induced via the homomorphism \(\pi_{1}(f,a)\) by an action of \(\pi_{1}(B,b)\) on \(E_{a}\) . This implies the injectivity of the homomorphism \(\pi_{1}(f,a)\) .

PROPOSITION 2. --- Let A be a connected and locally path-connected topological space, B a loop-disentanglable topological space, and let \(f\colon\mathrm{A}\to\mathrm{B}\) be a continuous map. Let a be a point of A and \(b=f(a)\) . The following assertions are equivalent:

\begin{enumerate}
\def\labelenumi{(\roman{enumi})}
\item
  The homomorphism \(\pi_1(f,a)\colon \pi_1(A,a)\to \pi_1(B,b)\) is surjective.
\item
  For every pair (E, E') of coverings of B, the map
\end{enumerate}

\[
f ^ {*} \colon \mathscr {C} _ {\mathrm{B}} (\mathrm{E}; \mathrm{E} ^ {\prime}) \rightarrow \mathscr {C} _ {\mathrm{A}} (\mathrm{A} \times_ {\mathrm{B}} \mathrm{E}; \mathrm{A} \times_ {\mathrm{B}} \mathrm{E} ^ {\prime})
\]

which, to a B-morphism \(g\colon \mathrm{E}\to \mathrm{E}^{\prime}\) , associates the A-morphism

\[
f ^ {*} (g) \colon \mathrm{A} \times_ {\mathrm{B}} \mathrm{E} \rightarrow \mathrm{A} \times_ {\mathrm{B}} \mathrm{E} ^ {\prime}, \quad (x, y) \mapsto (x, g (y))
\]

is bijective.

By means of the homomorphism \(\pi_1(f, a)\) , every \(\pi_1(B, b)\)-set is equipped with a structure of \(\pi_1(A, a)\) -set. Condition (i) is then equivalent to the following condition (i'):

(i') For every pair \((\mathrm{F},\mathrm{F}^{\prime})\) of \(\pi_{1}(\mathrm{B},b)\) -sets, every \(\pi_{1}(\mathrm{A},a)\) -morphism from F to \(F^{\prime}\) is a \(\pi_{1}(\mathrm{B},b)\) -morphism. Indeed, if the homomorphism \(\pi_{1}(f,a)\) is surjective, every \(\pi_{1}(\mathrm{A},a)\) -morphism of \(\pi_{1}(\mathrm{B},b)\) -sets is a \(\pi_{1}(\mathrm{B},b)\) -morphism. Conversely, take a \(\pi_{1}(\mathrm{B},b)\) -set F reduced to a point and set \(\mathrm{F}^{\prime}=\pi_{1}(\mathrm{B},b)/f_{*}(\pi_{1}(\mathrm{A},a))\) . The map from F to \(F^{\prime}\) whose image is \(f_{*}(\pi_{1}(\mathrm{A},a))\) is a \(\pi_{1}(\mathrm{A},a)\) -morphism but is not a \(\pi_{1}(\mathrm{B},b)\) -morphism if \(F^{\prime}\) is not reduced to a point, that is, if \(\pi_{1}(f,a)\) is not surjective.

Since the space B is loop-disentanglable, every \(\pi_{1}(B,b)\) -set is isomorphic to the \(\pi_{1}(B,b)\) -set \(E_{b}\) , where E is a covering of B (IV, p. 344, remark 1). The equivalence of (i') and (ii) therefore follows from prop. 2 of III, p. 310.

PROPOSITION 3. --- Let B be a loop-disentanglable topological space and let A be a connected and locally path-connected subspace of B (for example an open and connected subset). Let a be a point of A. The following properties are equivalent:

\begin{enumerate}
\def\labelenumi{(\roman{enumi})}
\item
  Every covering of B is trivializable over A.
\item
  The image of the homomorphism from \(\pi_{1}(\mathrm{A},a)\) into \(\pi_{1}(\mathrm{B},a)\) deduced from the canonical injection is reduced to the identity element.
\end{enumerate}

The implication (ii)⇒(i) follows from corollary 3 (III, p. 309).

Conversely, suppose that every covering of B is trivializable over A. This is in particular the case for the simply path-connected covering \((\lambda_{a}(\mathrm{B}),\varepsilon_{\mathrm{B}})\) (IV, p. 342, th. 1), so that the homomorphism \(\pi_{1}(\mathrm{A},a)\to\pi_{1}(\mathrm{B},a)\) is trivial (III, p. 309, cor. 3 of prop. 1).

